\section{必要な時間}

\subsection{今の見込み}

表 5 に，残りの計算の見込みを示す．時刻は 2026-11-01 15 時の状態である．GPU は 1 枚で，E76 と E77 が分け合っている．2 本より多く並べると GPU のメモリの合計が上限（24 GB）に近づき，計算が 5〜25 倍遅くなる（今朝に起きた）．このため E78 は E77 が終わってから 1 本で回す．

\begin{table}[h]
\centering\small
\caption{残りの計算の見込み（2026-11-01 15 時）}
\begin{tabularx}{\linewidth}{@{}p{0.08\linewidth}p{0.2\linewidth}p{0.2\linewidth}p{0.17\linewidth}L@{}}
\toprule
実験 & 内容 & 状態 & 残りの見込み & 見積もりの根拠 \\
\midrule
E76 & 筋 1〜3 の学び直し（GPU） & 42 ターゲットのうち 23 が終わった & 7〜9 時間（今夜に終わる） & 終わった 23 ターゲットの所要（1 つ 5〜63 分．\texttt{logs/e76.log}） \\
E77 & 比較の第 1 波（GPU） & 850 単位のうち 20 単位が終わった（AVT の世界の IEMOCAP 即興．6.8 時間） & 約 80 時間（3.3 日） & 終わった 20 単位から作った 1 単位の所要の式（下の 6.2 節） \\
E78 & 比較の第 2 波（GPU） & E77 の後に自動で始まる & 3〜11 日（不確か） & 検査役の試走から出した 173〜272 時間を，単独で回すと半分ほどになると見た値．試走は GPU が遅くなっていた時間帯（07〜08 時）に測ったので，大きめに出ている見込みである \\
E79 & 判定の解析（CPU） & E77 と E78 の後 & 2〜4 時間 & 検査役が人工の入力で測った値からの外挿 \\
\bottomrule
\end{tabularx}
\end{table}

今の設定のままだと，E77 が終わるのは 11 月 4〜5 日ごろ，E78 が終わるのは 11 月 7〜16 日ごろである．E79 の判定はその数時間後に出る．E78 の見積もりの幅が広いのは，E78 の 1 単位の所要を正常な GPU の状態で測れていないためである．E77 が終わって E78 が始まると，最初のターゲットの所要から見積もりを測り直せる．

\subsection{E77 の時間の内訳}

E77 の 1 単位は，1 つのターゲット（評価のターゲットか仮のターゲット）の 1 回の反復で，全手法の全設定の候補を解く計算である．終わった 20 単位の実測から，1 単位の所要は「11.4 分＋候補の発話 1,000 本ごとに 1.22 分」で表せた（AVT の世界．候補は IEMOCAP 即興で約 2,350 本，MOSEI で約 15,100 本，MELD で約 9,200 本）．この式は IEMOCAP 即興の 20 単位の合計 6.8 時間を再現する（式の値 6.75 時間）．

時間の大部分は仮のターゲットの計算である．IEMOCAP 即興では，評価の 10 反復に 2.2 時間，仮のターゲット（MOSEI と MELD を各 5 回）に 3.8 時間かかった．仮のターゲットでは，すべての手法のすべての設定の候補を解いて，その中から設定を選ぶからである．

手法ごとに見ると，共正則化が実時間 6.8 時間のうち 2.4 時間（35\%）を使っていた（表 6）．共正則化は 1 単位あたり約 7 分かかり，この時間は候補の数にほとんど依らない．

\begin{table}[h]
\centering\small
\caption{E77 の AVT の世界の IEMOCAP 即興での手法ごとの所要（\texttt{results/e77\_fuse-audio-video-text/IEMOCAP-impro\_conv.csv}．割合は実時間 6.8 時間（24,524 秒．\texttt{logs/e77b\_fuse-audio-video-text.log}）に対する値）}
\begin{tabular}{@{}lrrrr@{}}
\toprule
手法 & 評価（時間） & 仮のターゲット（時間） & 合計（時間） & 割合 \\
\midrule
共正則化 & 1.18 & 1.22 & 2.40 & 35\% \\
ours・oursx と除去実験の弱いラベルの行 & 0.15 & 0.65 & 0.80 & 12\% \\
ラベル伝播 & 0.10 & 0.43 & 0.53 & 8\% \\
エントロピー最小化 & 0.08 & 0.43 & 0.51 & 7\% \\
自己学習 & 0.11 & 0.36 & 0.47 & 7\% \\
共訓練 & 0.13 & 0.23 & 0.36 & 5\% \\
参照の行 & 0.29 & — & 0.29 & 4\% \\
温度つき蒸留 & 0.04 & 0.20 & 0.24 & 4\% \\
\bottomrule
\end{tabular}
\end{table}

\note{残りの約 1.2 時間は，記録の区分を持たない計算（情報最大化，T3A，共訓練の内側の解き直し）と，データの読み込みと保存である．}



\section{削減案}

\begin{tcolorbox}
\textbf{案}：E77 の残りと E78 で，結果の値を変えない省き方を 1 つ（下の 1），比較の判定に使わない行の省き方を 1 つ（下の 2）入れる．E77 の残りは約 80 時間から約 58 時間に減り，E77 は約 1 日早く終わる見込みである．E78 も 1 の分（1 割ほど）短くなる．
\end{tcolorbox}

\subsection{1．IEMOCAP 台本の仮のターゲットの計算を，即興の結果の写しにする}

\begin{itemize}
\item \textbf{何をするか}：IEMOCAP 台本を評価のターゲットにしたときの仮のターゲットの計算をやめ，IEMOCAP 即興のときに計算した同じ単位の結果を写して使う．
\item \textbf{理由}：仮のターゲットの計算に使うものは，即興と台本で全部同じである．即興と台本は同じ話者なので，どちらがターゲットでも，もう一方はソースから外れる．そのため使うソースの集合（AVT の世界なら MOSEI と MELD）が同じになる．出発点の 5 つの設定値も，5 種類の生徒すべてで一致することを確かめた．データの分け方も同じである．つまり今は同じ計算を 2 回している．説明資料 p5 でも，即興と台本で選ばれた設定が完全に同じだった．
\item \textbf{確かめ方}：今走っている E77 の台本のターゲットの最初の 1 単位（MOSEI の反復 0）が，即興の同じ単位の出力と一致することを確かめてから写す．
\item \textbf{減る時間}：E77 の残りで約 11 時間，E78 で 1 割ほど．
\item \textbf{失うもの}：無い．
\end{itemize}

\subsection{2．共正則化は仮のターゲットで解かず，評価のターゲットでは相手に有利な値で比べる}

\begin{itemize}
\item \textbf{何をするか}：共正則化だけ，仮のターゲットでは計算しない．評価のターゲットでは今どおり 4 つの設定を全部解き，テストのラベルで選んだ最良の値（オラクルの値）と oursx を比べる．
\item \textbf{理由}：共正則化は E77 の実時間の 35\% を使う．比較の設計では群 C の探索の行で，筋 4 と筋 5 の判定（P4，P5-1〜P5-4）には使わない．テストで選んだ最良の値は，仮のターゲットで選んだ値より必ず高いか等しい．このため oursx がオラクルの値を上回れば，仮のターゲットで選んだ値も上回ったと言える．相手を不利にする向きの偏りは入らない．
\item \textbf{減る時間}：E77 の残りで約 13 時間（AVT の世界で約 3.5 時間，AV の世界で約 9.4 時間）．
\item \textbf{失うもの}：共正則化の「仮のターゲットで選んだ値」と「選択の取りこぼし」の列．oursx がオラクルの値を上回らなかった単位では，仮のターゲットで選んだ値との比較について何も言えない．E61 では共正則化の UAR は基準とほぼ同じか下だった．
\end{itemize}

\subsection{進め方}

\begin{enumerate}
\item 台本を直す（1 と 2 を引数で選べるようにする）．E79 の設計の群 C の探索の行を「共正則化はオラクルの値」と書き直す．
\item 検査役（gate-code，gate-compare）に差分を検査させる．
\item 今走っている E77 は，台本のターゲットの仮のターゲットを計算している途中である．これを止めて，直した台本で台本のターゲットから流し直す．止めて失うのは，写しで置き換わる計算だけである．E78 の待ち受けはそのまま使える．
\end{enumerate}

1 と 2 の作業と検査に 1〜2 時間かかる．その間も E77 は今の設定のまま進む．
